\question{Can aggregated results be de-aggregated, or reconstructed from repeated analyses?}
\label{sec:aggdisclosure}
\statusMethod

\begin{shortanswer}
For the AI training path the exposure is bounded and addressed: what leaves a practice is a single
coefficient vector per round, protected by the mechanisms above. For the \emph{federated analytics}
path, the structural answers carry over unchanged --- but the disclosure-control questions do not,
because they depend on which indicators are computed, at what granularity, and how often. That
specification lives with another work package and we cannot answer for it. What we can say is that
``analytics'' confers no automatic anonymity: an aggregate reaching the server is aggregated
\emph{across patients within one practice}, not across practices, and it arrives attributable.
\end{shortanswer}

\subsection{The two paths are one mechanism}

The assessment treats the analytics path (aggregated indicator vectors) and the training path
(gradients) as if they had different technical properties, and assigns them different legal
prospects. It is worth recording precisely what does and does not follow.

\textbf{``Federated analytics'' is not an architectural boundary in the software we use.} The term
appears nowhere in Flower \FlwrVersion{}; the package offers one message API and one strategy
interface. The term itself is a real and standard one --- it originates with Google
\cite{fedanalytics} and describes computing aggregate statistics over decentralised data without
collecting it --- and it fairly describes the indicator path. But in implementation, both paths are
the same mechanism: a client returning a record to a server through the SuperLink.

Three consequences follow directly, and they are answers rather than deferrals:

\begin{itemize}
  \item \textbf{Who sees what is identical.} Aggregation is still the terminal step; messages still
        carry their sender's identity; the contribution is still practice-attributable. An indicator
        vector does not become anonymous by virtue of being called an aggregate --- it is aggregated
        over the patients \emph{within} one practice, which is exactly the population in which small
        cells become disclosive.
  \item \textbf{Secure aggregation applies identically, and more cleanly.} An indicator vector is a
        numeric vector, so it masks and sums without difficulty --- unlike the tree payload
        (\S\ref{sec:secagg}).
  \item \textbf{Differential privacy applies, and is \emph{easier} to reason about here.} A counting
        query has sensitivity exactly 1 per patient, so the accountant is tight and standard. On the
        training path the sensitivity is bounded only indirectly, by the clipping norm acting as a
        proxy. The analytics path is the more tractable of the two for formal guarantees, not the
        less.
\end{itemize}

\subsection{What genuinely differs --- and what it needs}

The difference is the payload's \emph{semantics}, not its transport. A coefficient is a smooth
function of hundreds of records; a count over a thin stratum can be one patient. Hence the controls
the CheCKD opinion asks for --- minimum cell counts, primary and secondary suppression,
$k$-anonymity --- which are meaningful for counts and largely meaningless for coefficients.

The specifically difficult case counsel raises is \textbf{repeated or overlapping analyses}. Two
queries differing by one predicate reveal the difference between them; a sequence of such queries can
isolate an individual even when no single result violates a cell-count floor. Suppression alone does
not defend against this, because each release passes its own test. Defending against it requires
treating the sequence as a whole: a query log, a cumulative disclosure budget spent across releases,
and refusal once it is exhausted.

\textbf{This test bed implements no such query auditing}, because it implements no analytics path.
Its metric-hygiene layer (\S\ref{sec:sidechannels}) applies a cell-count floor to per-practice metric
lines, which is the primitive but not the system.

\begin{sidenote}{For the analytics work package --- what we would need to answer this properly}
We can apply the same accountant and a cell-suppression analysis to the analytics path, but not
without its specification. Concretely:
\begin{enumerate}[label=(\alph*)]
  \item the list of indicators and their definitions;
  \item the stratification variables and the finest granularity any release may use;
  \item the query cadence --- how often indicators are recomputed and republished;
  \item whether the same cohort is re-queried over time as it changes, which is what makes
        differencing possible;
  \item whether any cumulative disclosure budget is tracked across releases today, and if so how.
\end{enumerate}
With (a)--(e) we can state the minimum cell count and the composed budget the analytics path needs.
Without them, this report answers only the architectural half, and counsel should read the
disclosure-control half as open rather than as settled.
\end{sidenote}
